\documentclass[11pt,a4paper]{article}
\usepackage{fontspec,amsmath,amssymb,amsthm,microtype}
\setmainfont{lmroman10-regular.otf}[BoldFont=lmroman10-bold.otf,ItalicFont=lmroman10-italic.otf,BoldItalicFont=lmroman10-bolditalic.otf]
\usepackage{xcolor}
\usepackage[most]{tcolorbox}
\usepackage{mathrsfs}
\providecommand{\lt}{<}
\providecommand{\gt}{>}
\setmonofont{DejaVuSansMono.ttf}[BoldFont=DejaVuSansMono-Bold.ttf,ItalicFont=DejaVuSansMono-Oblique.ttf,BoldItalicFont=DejaVuSansMono-BoldOblique.ttf,Scale=0.86]
\newfontfamily\leanfallback{STIXTwoMath-Regular.otf}[Scale=0.86]
\definecolor{leanrule}{gray}{0.62}
\definecolor{leanback}{gray}{0.962}
\definecolor{leantitle}{gray}{0.22}
\definecolor{leanmuted}{gray}{0.38}
\newtcolorbox{leanbox}[1]{enhanced,breakable,colback=leanback,colframe=leanrule,
  boxrule=0.4pt,arc=0.6mm,left=2.4mm,right=2.4mm,top=1.8mm,bottom=1.8mm,
  before skip=0.8em,after skip=1em,fonttitle=\small\bfseries,coltitle=leantitle,
  colbacktitle=leanback,titlerule=0.3pt,title={#1}}
\newlength{\leancharwidth}
\newenvironment{leancode}{\par\smallskip\begingroup\ttfamily\footnotesize\raggedright
  \settowidth{\leancharwidth}{M}\setlength{\parskip}{0pt}\setlength{\parindent}{0pt}
  \frenchspacing}{\par\endgroup\smallskip}
\newcommand{\leanline}[2]{\hangindent=\dimexpr#1\leancharwidth+4\leancharwidth\relax
  \hangafter=1\noindent\hspace*{#1\leancharwidth}#2\par}
\newcommand{\leansource}[1]{\par{\footnotesize\color{leanmuted}#1}\par}
\newcommand{\leanlabel}[1]{\par\smallskip\noindent{\small\bfseries #1}\par}
\usepackage[margin=25mm]{geometry}
\usepackage[unicode,hidelinks]{hyperref}
\hypersetup{pdfauthor={LangevinLean Project}}
\setlength{\emergencystretch}{3em}
\setlength{\parskip}{0.4em}
\setlength{\parindent}{0pt}
\allowdisplaybreaks[2]
\raggedbottom
\title{\texorpdfstring{Scalar recurrences for the normalized error}{Scalar recurrences for the normalized error}}
\author{LangevinLean Project}
\date{}
\begin{document}\maketitle
Draft mathematical exposition; these notes may contain mistakes.


This section solves the scalar recurrences that turn a one-step Wasserstein bound into a bound after many steps. No probability enters: the statements concern real sequences that are assumed to satisfy a one-step inequality. In the applications the sequence is the normalized error \(r_k=\frac mp\mathsf W_2^2(\nu_k,\pi)\) of the Langevin Monte Carlo chain, \(t=mh\) is the dimensionless step, and \(a=\kappa_{\rm av}+1\) with \(\kappa_{\rm av}=M_{\rm av}/m\). Multiplying the one-step bound of Theorem 1 by \(m/p\) and dropping its nonpositive last term gives


\[r_{k+1}\le\frac{1-t}{1+t}\,r_k+\frac{a\,t^2}{1+t},\]


because \(\frac mp\cdot\frac{(M_{\rm av}+m)h^2p}{1+mh}=\frac{(\kappa_{\rm av}+1)(mh)^2}{1+mh}\). We prove the elementary inequality \(\frac{1-t}{1+t}\le e^{-2t}\) on \([0,1]\), the iteration of an affine inequality whose fixed point is known, the resulting constant-step bound \(r_n\le e^{-2tn}r_0+\frac{at}2\) (the scalar content of the second claim of Theorem 1), and the induction step of the decreasing-step schedule of Proposition 1, in which \(t_k=1/(\kappa+k)\) and \(r_k\le a/(\kappa+k)\).



Earlier results used here: \cite[Section~3.15]{LangevinLean2026ULASC}.

\section{\texorpdfstring{The contraction factor is at most an exponential}{The contraction factor is at most an exponential}}


\textbf{Claim.} For every real \(t\) with \(0\le t\le1\),


\[\frac{1-t}{1+t}\le e^{-2t}.\]


\textbf{Proof.} If \(t=1\), the left-hand side is \(0\) and the right-hand side \(e^{-2}\) is positive. Let \(0\le t<1\), so that \(1+t>0\) and \(1-t>0\). Consider \(\varphi(s)=\log(1+s)-\log(1-s)-2s\) on \([0,1)\). Then \(\varphi(0)=0\) and


\[\varphi'(s)=\frac1{1+s}+\frac1{1-s}-2=\frac{2}{1-s^2}-2=\frac{2s^2}{1-s^2}\ge0,\]


so \(\varphi\) is nondecreasing by the mean value theorem and \(\varphi(t)\ge0\), that is,


\[2t\le\log\frac{1+t}{1-t}.\qquad(1)\]


(Equivalently, \(t\) is the first term of the series \(\frac12\log\frac{1+t}{1-t}=\sum_{i\ge0}\frac{t^{2i+1}}{2i+1}\), whose terms are nonnegative.) Since \(\frac{1-t}{1+t}=\bigl(\frac{1+t}{1-t}\bigr)^{-1}\) and \(\log(x^{-1})=-\log x\) for \(x>0\), (1) gives \(\log\frac{1-t}{1+t}\le-2t\). The exponential is increasing and \(\frac{1-t}{1+t}>0\), so exponentiating gives the claim. \(\square\)




In the applications \(t=mh\), and \(t\le1\) follows from \(m\le M_\infty\) and \(hM_\infty\le1\).

\begin{leanbox}{Lean: \texttt{contraction\_\allowbreak{}le\_\allowbreak{}exp}}
\leansource{\texttt{contraction\_\allowbreak{}le\_\allowbreak{}exp} \quad AverageLMC/\allowbreak{}Recurrence.lean, lines 21–34}
\begin{leancode}
\leanline{0}{theorem\ contraction\_le\_exp\ \{t\ :\ ℝ\}\ (ht\ :\ 0\ ≤\ t)\ (ht1\ :\ t\ ≤\ 1)\ :}
\leanline{4}{(1\ -\ t)\ /\ (1\ +\ t)\ ≤\ Real.exp\ (-2\ *\ t)}
\end{leancode}
\leanlabel{How the Lean proof follows the argument}
\texttt{contraction\_\allowbreak{}le\_\allowbreak{}exp} takes \texttt{ht : 0 \ensuremath{\le} t} and \texttt{ht1 : t \ensuremath{\le} 1}. \texttt{rcases eq\_\allowbreak{}or\_\allowbreak{}lt\_\allowbreak{}of\_\allowbreak{}le ht1} splits off the case \(t=1\), closed by \texttt{simpa} with the positivity of \(e^{-2}\) (\texttt{Real.\allowbreak{}exp\_\allowbreak{}pos}). For \(t<1\), \texttt{hl} is the Mathlib inequality \texttt{Real.\allowbreak{}sum\_\allowbreak{}range\_\allowbreak{}le\_\allowbreak{}log\_\allowbreak{}div ht h 1}, which states \(\sum_{i<n}t^{2i+1}/(2i+1)\le\frac12\log\frac{1+t}{1-t}\); with \(n=1\) and after \texttt{simp only} it is (1). \texttt{hi} writes the contraction as the inverse of \((1+t)/(1-t)\). The goal is moved to logarithms with \texttt{Real.\allowbreak{}log\_\allowbreak{}le\_\allowbreak{}iff\_\allowbreak{}le\_\allowbreak{}exp} (valid because the quotient is positive, \texttt{div\_\allowbreak{}pos hm hp}), \texttt{Real.\allowbreak{}log\_\allowbreak{}inv} turns the logarithm of the inverse into \(-\log\frac{1+t}{1-t}\), and \texttt{linarith} concludes from \texttt{hl}.
\end{leanbox}

\section{\texorpdfstring{Iterating an affine inequality over a finite horizon}{Iterating an affine inequality over a finite horizon}}


\textbf{Claim.} Let \((r_k)_{k\ge0}\) be real numbers, let \(q\ge0\), \(c\ge0\) and \(n\in\mathbb N\), and suppose that


\[r_{j+1}\le q\,r_j+(1-q)\,c\qquad\text{for every }j<n.\]


Then \(r_n\le q^nr_0+c\).




\textbf{Proof.} We first prove by induction on \(k\) that


\[r_k\le q^kr_0+(1-q^k)\,c\qquad\text{for every }k\le n.\qquad(2)\]


\emph{Base case \(k=0\).} Since \(q^0=1\), the right-hand side is \(r_0\).




\emph{Inductive step.} Let \(k+1\le n\) and assume (2) for \(k\) (which is allowed since \(k\le n\)). Because \(q\ge0\), multiplying the induction hypothesis by \(q\) preserves it: \(q\,r_k\le q^{k+1}r_0+(q-q^{k+1})c\). Since \(k<n\), the hypothesis applies with \(j=k\):


\[r_{k+1}\le q\,r_k+(1-q)c\le q^{k+1}r_0+(q-q^{k+1}+1-q)\,c=q^{k+1}r_0+(1-q^{k+1})\,c.\]


This proves (2). For \(k=n\), (2) gives \(r_n\le q^nr_0+c-q^nc\le q^nr_0+c\), because \(q^nc\ge0\) as a product of the nonnegative numbers \(q^n\) and \(c\). \(\square\)




The constant \(c\) is the fixed point of the map \(x\mapsto qx+(1-q)c\). The hypothesis is required only for the first \(n\) indices, which is what allows the bound to be applied to the warmup phase of a schedule that changes after \(n\) steps. No condition \(q\le1\) is needed.

\begin{leanbox}{Lean: \texttt{affine\_\allowbreak{}floor\_\allowbreak{}up\_\allowbreak{}to}}
\leansource{\texttt{affine\_\allowbreak{}floor\_\allowbreak{}up\_\allowbreak{}to} \quad AverageLMC/\allowbreak{}Recurrence.lean, lines 37–52}
\begin{leancode}
\leanline{0}{theorem\ affine\_floor\_up\_to\ (r\ :\ ℕ\ →\ ℝ)\ (q\ c\ :\ ℝ)\ (hq\ :\ 0\ ≤\ q)}
\leanline{4}{(hc\ :\ 0\ ≤\ c)\ (n\ :\ ℕ)}
\leanline{4}{(hstep\ :\ ∀\ j\ <\ n,\ r\ (j\ +\ 1)\ ≤\ q\ *\ r\ j\ +\ (1\ -\ q)\ *\ c)\ :}
\leanline{4}{r\ n\ ≤\ q\ \textasciicircum{}\ n\ *\ r\ 0\ +\ c}
\end{leancode}
\leanlabel{How the Lean proof follows the argument}
In \texttt{affine\_\allowbreak{}floor\_\allowbreak{}up\_\allowbreak{}to} the sequence is \texttt{r : \ensuremath{\mathbb N} \ensuremath{\to} \ensuremath{\mathbb R}} and the finite-horizon hypothesis is \texttt{hstep : \ensuremath{\forall} j < n, r (j + 1) \ensuremath{\le} q * r j + (1 -\allowbreak{} q) * c}. The strengthened statement (2) is \texttt{hstrong}, proved by \texttt{induction k with}: the \texttt{zero} case is \texttt{simp}; in the \texttt{succ} case \texttt{hprev} is the induction hypothesis (applicable because \texttt{Nat.\allowbreak{}le\_\allowbreak{}of\_\allowbreak{}succ\_\allowbreak{}le hk} gives \(k\le n\)), \texttt{hmul} multiplies it by \(q\ge0\) with \texttt{mul\_\allowbreak{}le\_\allowbreak{}mul\_\allowbreak{}of\_\allowbreak{}nonneg\_\allowbreak{}left}, \texttt{hs} is \texttt{hstep k} (with \texttt{Nat.\allowbreak{}lt\_\allowbreak{}of\_\allowbreak{}succ\_\allowbreak{}le hk} for \(k<n\)), \texttt{rw [pow\_\allowbreak{}succ]} writes \(q^{k+1}=q^kq\), and \texttt{nlinarith} performs the displayed computation. The final \texttt{nlinarith} uses \(q^nc\ge0\), supplied as \texttt{mul\_\allowbreak{}nonneg (pow\_\allowbreak{}nonneg hq n) hc}.
\end{leanbox}

\section{\texorpdfstring{The constant-step bound}{The constant-step bound}}


\textbf{Claim.} Let \((r_k)_{k\ge0}\) be real numbers and \(a,t\in\mathbb R\) with \(a\ge0\), \(0\le t\le1\) and \(r_0\ge0\). Let \(n\in\mathbb N\) and suppose that


\[r_{j+1}\le\frac{1-t}{1+t}\,r_j+\frac{a\,t^2}{1+t}\qquad\text{for every }j<n.\qquad(3)\]


Then


\[r_n\le e^{-2tn}\,r_0+\frac{a\,t}2.\qquad(4)\]


In particular, if (3) holds for every \(j\in\mathbb N\), then (4) holds for every \(n\in\mathbb N\).




\textbf{Proof.} Put \(q=\frac{1-t}{1+t}\) and \(c=\frac{at}2\). Since \(0\le t\le1\), we have \(1+t>0\) and \(1-t\ge0\), hence \(q\ge0\); also \(c\ge0\) because \(a,t\ge0\). The forcing term in (3) equals \((1-q)c\):


\[(1-q)\,c=\frac{(1+t)-(1-t)}{1+t}\cdot\frac{at}2=\frac{2t}{1+t}\cdot\frac{at}2=\frac{a\,t^2}{1+t}.\]


Thus (3) is the hypothesis of Step 2 with these \(q\) and \(c\), and Step 2 gives \(r_n\le q^nr_0+\frac{at}2\). By Step 1, \(0\le q\le e^{-2t}\); raising to the power \(n\) preserves inequalities between nonnegative numbers, so \(q^n\le(e^{-2t})^n=e^{-2tn}\). Multiplying by \(r_0\ge0\) gives \(q^nr_0\le e^{-2tn}r_0\), which proves (4). The second assertion applies the first to each \(n\), the hypothesis for all \(j\) containing the hypothesis for \(j<n\). \(\square\)




With \(r_k=\frac mp\mathsf W_2^2(\nu_k,\pi)\), \(t=mh\) and \(a=\kappa_{\rm av}+1\), multiplying (4) by \(p/m\) gives \(\mathsf W_2^2(\nu_n,\pi)\le e^{-2mnh}\mathsf W_2^2(\nu_0,\pi)+\frac{(M_{\rm av}+m)hp}{2m}\), since \(\frac pm\cdot\frac{(\kappa_{\rm av}+1)mh}2=\frac{(M_{\rm av}+m)hp}{2m}\). This is the second claim of Theorem 1.

\begin{leanbox}{Lean: \texttt{constant\_\allowbreak{}step\_\allowbreak{}normalized\_\allowbreak{}up\_\allowbreak{}to}, \texttt{constant\_\allowbreak{}step\_\allowbreak{}normalized}}
\leansource{\texttt{constant\_\allowbreak{}step\_\allowbreak{}normalized\_\allowbreak{}up\_\allowbreak{}to} \quad AverageLMC/\allowbreak{}Recurrence.lean, lines 56–75}
\begin{leancode}
\leanline{0}{theorem\ constant\_step\_normalized\_up\_to\ (r\ :\ ℕ\ →\ ℝ)\ (a\ t\ :\ ℝ)}
\leanline{4}{(ha\ :\ 0\ ≤\ a)\ (ht\ :\ 0\ ≤\ t)\ (ht1\ :\ t\ ≤\ 1)\ (hr0\ :\ 0\ ≤\ r\ 0)\ (n\ :\ ℕ)}
\leanline{4}{(hstep\ :\ ∀\ j\ <\ n,\ r\ (j\ +\ 1)\ ≤}
\leanline{6}{(1\ -\ t)\ /\ (1\ +\ t)\ *\ r\ j\ +\ a\ *\ t\ \textasciicircum{}\ 2\ /\ (1\ +\ t))\ :}
\leanline{4}{r\ n\ ≤\ Real.exp\ (-2\ *\ t\ *\ (n\ :\ ℝ))\ *\ r\ 0\ +\ a\ *\ t\ /\ 2}
\end{leancode}
\leansource{\texttt{constant\_\allowbreak{}step\_\allowbreak{}normalized} \quad AverageLMC/\allowbreak{}Recurrence.lean, lines 78–83}
\begin{leancode}
\leanline{0}{theorem\ constant\_step\_normalized\ (r\ :\ ℕ\ →\ ℝ)\ (a\ t\ :\ ℝ)}
\leanline{4}{(ha\ :\ 0\ ≤\ a)\ (ht\ :\ 0\ ≤\ t)\ (ht1\ :\ t\ ≤\ 1)\ (hr0\ :\ 0\ ≤\ r\ 0)}
\leanline{4}{(hstep\ :\ ∀\ j,\ r\ (j\ +\ 1)\ ≤}
\leanline{6}{(1\ -\ t)\ /\ (1\ +\ t)\ *\ r\ j\ +\ a\ *\ t\ \textasciicircum{}\ 2\ /\ (1\ +\ t))\ (n\ :\ ℕ)\ :}
\leanline{4}{r\ n\ ≤\ Real.exp\ (-2\ *\ t\ *\ (n\ :\ ℝ))\ *\ r\ 0\ +\ a\ *\ t\ /\ 2}
\end{leancode}
\leanlabel{How the Lean proof follows the argument}
\begin{enumerate}


\item \texttt{constant\_\allowbreak{}step\_\allowbreak{}normalized\_\allowbreak{}up\_\allowbreak{}to} has the hypotheses \texttt{ha : 0 \ensuremath{\le} a}, \texttt{ht : 0 \ensuremath{\le} t}, \texttt{ht1 : t \ensuremath{\le} 1}, \texttt{hr0 : 0 \ensuremath{\le} r 0} and the finite-horizon recurrence \texttt{hstep}, which is (3). \texttt{hq} is \(q\ge0\) (\texttt{div\_\allowbreak{}nonneg}), and \texttt{hfix} is the identity \(\frac{at^2}{1+t}=(1-q)\frac{at}2\), proved by \texttt{field\_\allowbreak{}simp} and \texttt{ring}.

\item \texttt{h} applies \texttt{affine\_\allowbreak{}floor\_\allowbreak{}up\_\allowbreak{}to} (Step 2) with \(c=at/2\), after \texttt{simpa only [hfix]} rewrites the forcing term of \texttt{hstep}.

\item \texttt{hp} raises the inequality of \texttt{contraction\_\allowbreak{}le\_\allowbreak{}exp} to the \(n\)th power with \texttt{pow\_\allowbreak{}le\_\allowbreak{}pow\_\allowbreak{}left\ensuremath{{}_{0}}}, and \texttt{heq} identifies \((e^{-2t})^n\) with \(e^{-2tn}\) through \texttt{Real.\allowbreak{}exp\_\allowbreak{}nat\_\allowbreak{}mul}. The last line multiplies by \(r_0\ge0\) (\texttt{mul\_\allowbreak{}le\_\allowbreak{}mul\_\allowbreak{}of\_\allowbreak{}nonneg\_\allowbreak{}right hp hr0}) and chains with \texttt{h}.

\item \texttt{constant\_\allowbreak{}step\_\allowbreak{}normalized} is the version for a recurrence valid at every index: it calls the previous theorem with \texttt{fun j \_\allowbreak{} \ensuremath{\mapsto} hstep j}.

\end{enumerate}
\end{leanbox}

\section{\texorpdfstring{The induction step of the decreasing-step schedule}{The induction step of the decreasing-step schedule}}


\textbf{Claim.} (1) Let \(a,x,r\in\mathbb R\) with \(x\ge1\) and \(r\le a/x\). Then


\[\frac{1-1/x}{1+1/x}\,r+\frac{a\,(1/x)^2}{1+1/x}\le\frac a{x+1}.\]


(2) Let \((r_n)_{n\ge0}\) be real numbers and \(a,\kappa\in\mathbb R\) with \(a\ge0\), \(\kappa\ge1\) and \(r_0\le a/\kappa\). Suppose that for every \(n\in\mathbb N\), with \(t_n=1/(\kappa+n)\),


\[r_{n+1}\le\frac{1-t_n}{1+t_n}\,r_n+\frac{a\,t_n^2}{1+t_n}.\qquad(5)\]


Then \(r_n\le\frac a{\kappa+n}\) for every \(n\in\mathbb N\).




\textbf{Proof.} (1) Since \(x\ge1>0\), we have \(1/x\le1\) and \(1+1/x>0\), so the coefficient \(q=\frac{1-1/x}{1+1/x}\) is nonnegative. Multiplying \(r\le a/x\) by \(q\ge0\),


\[q\,r+\frac{a(1/x)^2}{1+1/x}\le q\,\frac ax+\frac{a(1/x)^2}{1+1/x}.\]


Multiplying numerators and denominators by \(x\) gives \(q=\frac{x-1}{x+1}\) and \(\frac{(1/x)^2}{1+1/x}=\frac1{x(x+1)}\), so the right-hand side equals


\[\frac{a(x-1)}{x(x+1)}+\frac a{x(x+1)}=\frac{a\,x}{x(x+1)}=\frac a{x+1}.\]


(2) Induction on \(n\). For \(n=0\) the claim is the hypothesis \(r_0\le a/\kappa\). Assume \(r_n\le a/(\kappa+n)\) and put \(x=\kappa+n\). Then \(x\ge1\) because \(\kappa\ge1\) and \(n\ge0\), and \(t_n=1/x\). By (5) and part (1),


\[r_{n+1}\le\frac{1-1/x}{1+1/x}\,r_n+\frac{a(1/x)^2}{1+1/x}\le\frac a{x+1}=\frac a{\kappa+(n+1)}.\qquad\square\]


Part (1) uses no sign condition on \(a\), and the hypothesis \(a\ge0\) of part (2) is not used in its proof. In Proposition 1, \(x=\kappa_\infty+k-k_0\) and \(t_n=mh_k\) is the decreasing step after the warmup phase; part (1) states that the bound \(a/x\) is propagated to \(a/(x+1)\) by one step of the schedule.

\begin{leanbox}{Lean: \texttt{harmonic\_\allowbreak{}step}, \texttt{variable\_\allowbreak{}step\_\allowbreak{}tail}}
\leansource{\texttt{harmonic\_\allowbreak{}step} \quad AverageLMC/\allowbreak{}Recurrence.lean, lines 86–100}
\begin{leancode}
\leanline{0}{theorem\ harmonic\_step\ \{a\ x\ r\ :\ ℝ\}\ (hx\ :\ 1\ ≤\ x)}
\leanline{4}{(hr\ :\ r\ ≤\ a\ /\ x)\ :}
\leanline{4}{(1\ -\ 1\ /\ x)\ /\ (1\ +\ 1\ /\ x)\ *\ r\ +\ a\ *\ (1\ /\ x)\ \textasciicircum{}\ 2\ /\ (1\ +\ 1\ /\ x)}
\leanline{6}{≤\ a\ /\ (x\ +\ 1)}
\end{leancode}
\leansource{\texttt{variable\_\allowbreak{}step\_\allowbreak{}tail} \quad AverageLMC/\allowbreak{}Recurrence.lean, lines 104–117}
\begin{leancode}
\leanline{0}{theorem\ variable\_step\_tail\ (r\ :\ ℕ\ →\ ℝ)\ (a\ κ\ :\ ℝ)}
\leanline{4}{(\_ha\ :\ 0\ ≤\ a)\ (hκ\ :\ 1\ ≤\ κ)\ (hinit\ :\ r\ 0\ ≤\ a\ /\ κ)}
\leanline{4}{(hstep\ :\ ∀\ n,\ r\ (n\ +\ 1)\ ≤}
\leanline{6}{(1\ -\ 1\ /\ (κ\ +\ (n\ :\ ℝ)))\ /\ (1\ +\ 1\ /\ (κ\ +\ (n\ :\ ℝ)))\ *\ r\ n\ +}
\leanline{6}{a\ *\ (1\ /\ (κ\ +\ (n\ :\ ℝ)))\ \textasciicircum{}\ 2\ /\ (1\ +\ 1\ /\ (κ\ +\ (n\ :\ ℝ))))}
\leanline{4}{(n\ :\ ℕ)\ :\ r\ n\ ≤\ a\ /\ (κ\ +\ (n\ :\ ℝ))}
\end{leancode}
\leanlabel{How the Lean proof follows the argument}
\begin{enumerate}


\item \texttt{harmonic\_\allowbreak{}step} has the hypotheses \texttt{hx : 1 \ensuremath{\le} x} and \texttt{hr : r \ensuremath{\le} a /\allowbreak{} x}. The auxiliary facts are \(x>0\) (\texttt{hx0}), \(1+1/x>0\) (\texttt{hd}), \(1/x\le1\) (\texttt{hdiv}, from \texttt{div\_\allowbreak{}le\_\allowbreak{}one}) and \(q\ge0\) (\texttt{hq}). The \texttt{calc} block has two lines: the first multiplies \texttt{hr} by \(q\) with \texttt{mul\_\allowbreak{}le\_\allowbreak{}mul\_\allowbreak{}of\_\allowbreak{}nonneg\_\allowbreak{}left} and adds the unchanged forcing term (\texttt{add\_\allowbreak{}le\_\allowbreak{}add … le\_\allowbreak{}rfl}); the second is the algebraic identity, proved by \texttt{field\_\allowbreak{}simp} and \texttt{ring}.

\item \texttt{variable\_\allowbreak{}step\_\allowbreak{}tail} has the hypotheses \texttt{\_\allowbreak{}ha : 0 \ensuremath{\le} a} (unused, as the leading underscore records), \texttt{h\ensuremath{\kappa} : 1 \ensuremath{\le} \ensuremath{\kappa}}, \texttt{hinit : r 0 \ensuremath{\le} a /\allowbreak{} \ensuremath{\kappa}}, and \texttt{hstep}, which is (5) written with \(t_n\) spelled out as \texttt{1 /\allowbreak{} (\ensuremath{\kappa} + (n : \ensuremath{\mathbb R}))}. The proof is \texttt{induction n with}: the \texttt{zero} case is \texttt{simpa using hinit}; in the \texttt{succ} case \texttt{hx} proves \(\kappa+n\ge1\) from \texttt{Nat.\allowbreak{}cast\_\allowbreak{}nonneg}, \texttt{(hstep n).\allowbreak{}trans (harmonic\_\allowbreak{}step hx ih)} chains (5) with part (1), and \texttt{simpa only [Nat.\allowbreak{}cast\_\allowbreak{}add, Nat.\allowbreak{}cast\_\allowbreak{}one, add\_\allowbreak{}assoc]} rewrites \(\kappa+n+1\) as \(\kappa+(n+1)\).

\end{enumerate}
\end{leanbox}

\begin{thebibliography}{ULA-LSI}
\bibitem[ULA-SC]{LangevinLean2026ULASC}
LangevinLean Project. \emph{Unadjusted Langevin under Strong Convexity}. LL-01, 2026. Draft report, content revision d50c0f64815b. \url{https://odakyildiz.github.io/LangevinLean/docs/pdf/UnadjustedLangevin/FullProof.pdf}.
\end{thebibliography}

\end{document}
